% =============================================================================
% KAPITEL 5 – Evaluation und Auswahl des Observability-Stacks
% =============================================================================


\chapter{Evaluation und Auswahl des Observability-Stacks}
 
\section{Festlegung der Vergleichskriterien}
 
Die acht Vergleichskriterien für die Stack-Auswahl wurden in Abschnitt~4.3 aus den Anforderungen A1--A5 und der einschlägigen Literatur hergeleitet. Sie umfassen die Abdeckung der drei Telemetriepfeiler, Korrelationsfähigkeit, OpenTelemetry-Kompatibilität, Dashboarding und Alerting, Integrationsaufwand, Ressourcenbedarf, Erweiterbarkeit sowie Community und Ökosystem. Im Kern adressieren die Anforderungen drei Kernbedürfnisse, nämlich eine durchgängige Nachvollziehbarkeit dienstübergreifender Abläufe (A1), eine mehrstufige Messbarkeit vom Infrastruktur- bis zum Anwendungsniveau (A2) sowie ein korrelierbares, strukturiertes Logging (A3). Hinzu kommen die Betriebsrestriktionen einer Single-Host-Topologie (A4) und das Erfordernis, alle Stakeholder-Rollen mit handlungsorientierten Informationen zu versorgen (A5). Das vorliegende Kapitel wendet diese Kriterien auf drei vollständige Observability-Architekturen an, um Forschungsfrage~2 zu beantworten: Welche Observability-Stack-Kombination eignet sich am besten für die Überwachungsanforderungen von frag.jetzt?
 
 
\section{OpenTelemetry als verbindliche Instrumentierungsschicht}
 
Bevor ein Vergleich der Analyse- und Visualisierungs-Backends sinnvoll ist, muss die vorgelagerte Instrumentierungsschicht festgelegt werden. Bereits in Kapitel~2 wurde OpenTelemetry als herstellerneutraler CNCF-Standard eingeordnet, der die Erzeugung und Erfassung von Logs, Metriken und Traces vereinheitlicht, ohne selbst Speicherung oder Darstellung zu übernehmen \parencite[S.~72030]{faseeha2025observability}. Für die vorliegende Arbeit wird OpenTelemetry als verbindliche Instrumentierungsgrundlage festgelegt. Die Begründung folgt aus drei Überlegungen.
 
Erstens entkoppelt eine standardisierte Instrumentierung die Entscheidung über den Applikationscode von der Entscheidung über das Backend. \textcite[S.~107]{sakinala2025observability} beschreiben den OpenTelemetry Collector als flexiblen, herstellerunabhängigen Verarbeitungsbaustein, der Telemetriedaten aus verschiedenen Quellen entgegennimmt, transformiert und an beliebige Zielsysteme weiterleitet. Sollte sich ein Backend künftig als ungeeignet erweisen, lässt sich die Analyse-Ebene austauschen, ohne die Instrumentierung im Anwendungscode anzupassen. Diese Austauschbarkeit mindert das Risiko einer Herstellerbindung. Ein Vorteil, den \textcite[S.~107]{sakinala2025observability} der herstellerneutralen Architektur von OpenTelemetry explizit zuschreiben.
 
Zweitens profitiert frag.jetzt als polyglotte Anwendung besonders von einem sprachübergreifenden Instrumentierungsansatz. Backend und WS-Gateway nutzen Java beziehungsweise Kotlin, der AI Provider basiert auf Python, und das Frontend ist in TypeScript geschrieben. OpenTelemetry stellt für alle vier Laufzeitumgebungen SDKs bereit, die einheitliche Semantikkonventionen und Datenmodelle verwenden \parencite[S.~106]{sakinala2025observability}. Ohne einen solchen gemeinsamen Rahmen müssten für jede Sprache separate Bibliotheken gepflegt werden. Ein Aufwand, der bei begrenzten Betriebsressourcen schnell unwirtschaftlich wird. \textcite[S.~15]{giamattei2024monitoring} empfehlen, Werkzeuge zu bevorzugen, deren Instrumentierung zum Entwicklungstempo des Projekts passt; die automatischen Instrumentierungsbibliotheken, die OpenTelemetry für gängige Frameworks bereitstellt, kommen dieser Anforderung entgegen.
 
Drittens sichert die Trägerschaft durch die Cloud Native Computing Foundation eine langfristige Weiterentwicklung und breite Herstellerunterstützung \parencite[S.~107]{sakinala2025observability}. Der folgende Vergleich setzt OpenTelemetry daher als gegebene Instrumentierungsschicht voraus und bewertet die Kandidaten-Stacks ausschließlich in ihrer Rolle als Empfangs-, Speicher-, Visualisierungs- und Alarmierungsebene.
 
 
\section{Vergleich der Zielarchitekturen}
 
Einzelne Werkzeuge wie Prometheus allein oder der klassische ELK-Stack ohne APM decken jeweils nur einen der drei Telemetriepfeiler vollwertig ab und kommen daher als alleinstehende Endkandidaten für Forschungsfrage~2 nicht in Betracht, was auch die Fachliteratur bestätigt. \textcite[S.~1, 5]{giamattei2024monitoring} führen Prometheus, Jaeger und Elasticsearch als komplementäre Werkzeuge, die typischerweise kombiniert werden; \textcite[S.~158--159]{garimilla2024designing} ordnen diese Werkzeuge explizit den Rollen Metrikerfassung, Tracing und Loganalyse zu, während Grafana die übergreifende Visualisierung übernimmt. Für den Vergleich werden deshalb ausschließlich Architekturen herangezogen, die Logs, Metriken und Traces vollständig adressieren und über eine gemeinsame Auswertungsebene verfügen. Alle drei Kandidaten basieren auf quelloffenen Kernkomponenten und werden in der Fachliteratur zur Microservices-Beobachtbarkeit regelmäßig herangezogen \parencite[S.~72019--72021]{faseeha2025observability}.
 
 
\subsection{Elastic Observability}
 
Elastic Observability erweitert den klassischen ELK-Stack um die APM-Komponente und OTel-Integration und adressiert damit alle drei Telemetriepfeiler innerhalb eines Herstellerökosystems. Der APM-Server unterstützt den Empfang von Traces, Metriken und Logs über das OpenTelemetry-Protokoll (OTLP) nativ \parencite{elastic_otel_docs}. Elasticsearch bildet als verteilte Such- und Analysemaschine den Kern und indiziert sämtliche eingehenden Dokumente im Volltext \parencite[S.~918]{banerjee2024comparative}. Kibana stellt die Visualisierungs- und Analyseschnittstelle bereit; die Trace-Log-Korrelation erfolgt innerhalb der Kibana-Oberfläche, während Metrik-Auswertungen über Lens-Visualisierungen angebunden werden \parencite{elastic_otel_docs}.
 
Die Stärke des Ansatzes liegt in der tiefgreifenden Loganalyse. Unstrukturierte Textdaten lassen sich mit hoher Geschwindigkeit durchsuchen, filtern und korrelieren \parencite[S.~920]{banerjee2024comparative}. Für den Kontext von frag.jetzt ergeben sich allerdings gewichtige Nachteile. Der Ansatz, sämtliche Logzeilen vollständig zu indizieren, verursacht erheblichen Speicher- und Rechenaufwand \parencite[S.~918]{banerjee2024comparative}. Auf der Single-Host-Topologie von frag.jetzt konkurriert ein Elasticsearch-Knoten mit seinem JVM-Heap unmittelbar mit den Anwendungscontainern um CPU und Arbeitsspeicher. Hinzu kommt die Lizenzentwicklung. Elasticsearch und Kibana wurden 2021 von Apache~2.0 auf die Elastic License~v2 und SSPL umgestellt; seit 2024 steht der Quellcode zusätzlich unter AGPL zur Verfügung \parencite{elastic_license_docs}. Aus Sicht des Autors schränkt die mehrfache Lizenzänderung die langfristige Planungssicherheit im Vergleich zu CNCF-getragenen Projekten ein, auch wenn die aktuelle AGPL-Option den Zugang wieder erleichtert.
 
 
\subsection{Grafana-Stack: Prometheus + Loki + Tempo + Grafana}
 
Die zweite Zielarchitektur kombiniert Prometheus für Metriken, Loki für Logs und Tempo für Distributed Traces mit Grafana als gemeinsamer Visualisierungs- und Alarmierungsplattform. \textcite[S.~501]{gudimetla2025transitioning} beschreiben eine vergleichbare Referenzarchitektur, in der ein OpenTelemetry Collector als zentrale Routing-Schicht Telemetriedaten an Prometheus, Loki und einen Trace-Speicher weiterleitet, während Grafana die integrierte Visualisierung aller drei Datenquellen übernimmt.
 
Loki verfolgt bei der Log-Aggregation einen ressourcenschonenden Ansatz. Statt den vollständigen Inhalt jeder Logzeile zu indizieren, indiziert Loki ausschließlich Labels, Metadaten wie Dienstname, Umgebung oder Log-Level \parencite[S.~72019--72020]{faseeha2025observability}. Dieser Mechanismus, konzeptionell an das Label-Modell von Prometheus angelehnt \parencite[S.~86911]{usman2022survey}, reduziert den Speicher- und Verarbeitungsbedarf erheblich. \textcite[S.~502]{gudimetla2025transitioning} beziffern die Speicherersparnis auf 50--80\,\% gegenüber volltextbasierter Indizierung, weisen zugleich aber darauf hin, dass die Abfrageflexibilität eingeschränkt ist. Suchanfragen ohne Label-Eingrenzung durchlaufen sämtliche zugehörigen Chunks und sind entsprechend langsamer. Die Grafana-Loki-Dokumentation ergänzt, dass zu viele dynamische Labels die Zahl der Streams erhöhen und zu erheblichen Leistungseinbußen führen können \parencite{grafana_loki_docs}. Für frag.jetzt, dessen Logvolumen moderat ausfällt und dessen Label-Struktur durch die begrenzte Zahl von Diensten überschaubar bleibt, überwiegen die Vorteile gegenüber der Vollindizierung.
 
Tempo empfängt Trace-Daten über standardisierte Protokolle, darunter OTLP sowie Jaeger- und Zipkin-Formate \parencite[S.~8]{tzanettis2022data}. Die Grafana-Tempo-Dokumentation beschreibt Tempo als Backend, das Trace-Daten speichert und deren Verknüpfung mit Logs und Metriken innerhalb von Grafana ermöglicht \parencite{grafana_tempo_docs}. \textcite[S.~16]{sundstrom2025finding} bestätigen diese Architektur in einer prototypischen Implementierung und beschreiben Grafana als Oberfläche, die alle Observability-Signale vereint und die Navigation zwischen Metriken, Traces und Logs erlaubt.
 
Prometheus übernimmt die Metrikerfassung und -speicherung. Auch ohne ein zusätzliches Langzeitspeicher-Backend wie Mimir ergibt sich eine zentrale Auswertungsoberfläche für alle drei Signaltypen, da Grafana Prometheus-, Loki- und Tempo-Datenquellen nativ unterstützt \parencite[S.~16--17]{sundstrom2025finding}. Mimir ist daher nicht Voraussetzung für eine einheitliche Oberfläche, sondern ein Prometheus-kompatibles, skalierbares Langzeitspeicher- und Abfragebackend \parencite{grafana_mimir_docs}. \textcite[S.~502]{gudimetla2025transitioning} beschreiben Prometheus als Lösung für Datenerfassung und Kurzzeitspeicherung, während Mimir, Thanos oder Cortex die horizontale Skalierung und Langzeitvorhaltung übernehmen. Für die aktuelle Skalierungsstufe von frag.jetzt ist die lokale Prometheus-TSDB für den Prototyp ausreichend. Mimir bleibt als evolutive Ausbauoption bestehen.
 
 
\subsection{Prometheus + Jaeger + Loki + Grafana}
 
Die dritte Zielarchitektur verfolgt einen Bausteinansatz aus jeweils eigenständigen, ausgereiften Open-Source-Werkzeugen. Prometheus übernimmt die Metriken, Jaeger das Distributed Tracing und Loki die Log-Aggregation, visualisiert über Grafana. Alle drei Backend-Komponenten lassen sich in eine OpenTelemetry-basierte Pipeline integrieren. Der OTel Collector leitet die jeweiligen Telemetriedaten an das zuständige Backend weiter \parencite[S.~501]{gudimetla2025transitioning}. Prometheus und Jaeger sind graduated CNCF-Projekte. \textcite[S.~15]{giamattei2024monitoring} heben die aktive Open-Source-Community beider Werkzeuge als langfristigen Stabilitätsfaktor hervor. Jaeger bietet eine dedizierte Trace-Oberfläche, die in der Literatur häufig als Referenz herangezogen wird \parencite[S.~72020]{faseeha2025observability}. \textcite[S.~72019]{faseeha2025observability} führen Prometheus, Loki und Jaeger als verbreitete Werkzeuge für Microservices-Observability auf, was zeigt, dass diese Kombination keine exotische Eigenkonstruktion darstellt.
 
Dem Vorteil der Komponentenreife steht eine erhöhte Integrationskomplexität gegenüber. Die drei Werkzeuge folgen unterschiedlichen Konfigurationsformaten und Storage-Modellen. Prometheus nutzt eine lokale TSDB, Jaeger benötigt ein separates Speicher-Backend wie Elasticsearch, Cassandra oder Badger \parencite[S.~501]{gudimetla2025transitioning}, und Loki folgt einem eigenen Chunk-Storage-Modell. Die Korrelation zwischen den Telemetriearten ist über Grafana grundsätzlich möglich. \textcite[S.~17]{sundstrom2025finding} weisen jedoch darauf hin, dass Jaeger und Zipkin typischerweise als separate Trace-Backends dienen, die über zusätzliche Konfiguration in Grafana eingebunden werden. Die Verknüpfung bleibt damit nach Einschätzung des Autors weniger eng integriert als in einer Architektur, deren Komponenten von vornherein aufeinander abgestimmt sind.
 
 
\section{Vergleich und Synthese}
 
Tabelle~\ref{tab:stackvergleich} fasst die Bewertung der drei Zielarchitekturen entlang der acht Vergleichskriterien zusammen. Die Skala reicht von 1 (eingeschränkt) über 2 (befriedigend) bis 3 (sehr gut). Die Einstufungen beruhen auf der qualitativen Analyse in den vorherigen Abschnitten und den in Kapitel~4.3 formulierten Kriterienausprägungen. Sie stellen eine argumentativ begründete Einschätzung des Autors dar, keine empirisch erhobenen Messwerte.

\begin{table}[ht]
\centering
\caption[Vergleichende Bewertung der drei Observability-Zielarchitekturen]{Vergleichende Bewertung der drei Observability-Zielarchitekturen (3\,=\,sehr gut, 2\,=\,befriedigend, 1\,=\,eingeschränkt). Alle drei Kandidaten adressieren Logs, Metriken und Traces. OpenTelemetry wird als gemeinsame Instrumentierungsschicht vorausgesetzt. \small Quelle: Eigene Darstellung.}
\label{tab:stackvergleich}
\small
\setlength{\tabcolsep}{3pt}
\renewcommand{\arraystretch}{1.1}
\begin{tabularx}{\textwidth}{
  >{\RaggedRight\arraybackslash}p{2.0cm}
  S >{\RaggedRight\arraybackslash}X
  S >{\RaggedRight\arraybackslash}X
  S >{\RaggedRight\arraybackslash}X
}
\toprule
\textbf{Kriterium}
  & \multicolumn{2}{c}{\textbf{Elastic Observability}}
  & \multicolumn{2}{c}{\textbf{Grafana-Stack}}
  & \multicolumn{2}{c}{\textbf{Prom./Jaeger/Loki}} \\
  & \multicolumn{2}{c}{\scriptsize (Elastic APM, ES, Kibana)}
  & \multicolumn{2}{c}{\scriptsize (Prom., Loki, Tempo, Grafana)}
  & \multicolumn{2}{c}{\scriptsize (Prom., Jaeger, Loki, Grafana)} \\
\midrule
Telemetrie\-pfeiler
  & 3 & Logs, Metriken und Traces über APM/OTel abgedeckt
  & 3 & Alle drei Pfeiler durch Prometheus, Loki und Tempo
  & 3 & Alle drei Pfeiler durch Einzelkomponenten \\[4pt]
Korrelation
  & 2 & Trace-Log-Korrelation in Kibana; Metrik-Anbindung über Lens
  & 3 & Trace-Log-Metrik-Navigation nativ in Grafana
  & 2 & Via Grafana-Plugins, weniger eng integriert \\[4pt]
OTel-Kompa\-tibilität
  & 3 & OTLP-Empfang nativ, EDOT-Distributionen verfügbar
  & 3 & Alle Komponenten OTel-nativ
  & 3 & Alle Komponenten OTel-nativ \\[4pt]
Dashboard.\ \& Alerting
  & 2 & Kibana + Watcher/Rules; weniger flexibel als Grafana
  & 3 & Grafana Unified Alerting + Alertmanager
  & 3 & Grafana + Alertmanager \\[4pt]
Integrations\-aufwand
  & 1 & Logstash/Beats-Pipeline, APM-Server, JVM-Tuning
  & 2 & Mehr Container, aber homogene Konfiguration via OTel-Collector
  & 1 & Drei Konfig.-Formate und Storage-Backends \\[4pt]
Ressourcen\-bedarf
  & 1 & Elasticsearch speicherintensiv, JVM-Heap konkurriert mit Anwendung
  & 2 & Höher als Prom.\ allein, deutlich geringer als Elastic
  & 2 & Jaeger benötigt separates Speicher-Backend \\[4pt]
Erweiter\-barkeit
  & 2 & Nativ skalierbar, aber Cluster-Management komplex
  & 3 & Prometheus via Mimir erweiterbar; Loki/Tempo modular
  & 3 & Jede Komponente einzeln skalierbar \\[4pt]
Community \& Ökosystem
  & 2 & Große Community, aber Lizenzänderungen (ELv2/SSPL, seit 2024 teils AGPL)
  & 2 & Grafana Labs getragen, wachsend, Tempo jünger als Jaeger
  & 3 & Prom.\ u.\ Jaeger graduated CNCF, Loki etabliert \\
\midrule
\textbf{Summe}
  & \multicolumn{2}{c}{\textbf{16}}
  & \multicolumn{2}{c}{\textbf{21}}
  & \multicolumn{2}{c}{\textbf{20}} \\
\bottomrule
\end{tabularx}
\end{table}

Bei den Telemetriepfeilern und der OTel-Kompatibilität schneiden alle drei Kandidaten gleichwertig ab, was eine direkte Folge der Konstruktionsentscheidung ist, nur vollständige Architekturen zu vergleichen. Die Differenzierung ergibt sich in den übrigen Kriterien. Elastic Observability erzielt die niedrigste Gesamtbewertung, weil der hohe Ressourcenbedarf von Elasticsearch und der komplexe Integrationspfad auf einer Single-Host-Topologie besonders ins Gewicht fallen. Die beiden Grafana-basierten Architekturen liegen eng beieinander. Der Unterschied manifestiert sich in den Kriterien Korrelation und Integrationsaufwand.
 
Beide Grafana-basierten Architekturen ermöglichen eine durchgängige Sicht auf das Laufzeitverhalten, indem sie anwendungsbezogene Entwurfsmuster und infrastrukturnahe Messgrößen als komplementäre Perspektiven verbinden \parencite[S.~22]{albuquerque2025tracing}. \textcite[S.~57]{bissig2024unified} bestätigen, dass eine einheitliche, auf OpenTelemetry und Distributed Tracing gestützte Analyseperspektive in heterogenen verteilten Systemen die Voraussetzung für umfassende Observability darstellt. Die Entscheidung zwischen den beiden Grafana-basierten Varianten wird im folgenden Abschnitt anhand des konkreten Betriebskontexts begründet.
 
 
\section{Begründete Zielentscheidung für frag.jetzt}
 
\subsection{Entscheidungsbegründung im Kontext von frag.jetzt}
 
Das erste und schwerwiegendste Argument betrifft die \emph{Integrationskomplexität bei begrenzten Betriebsressourcen}. \textcite[S.~9]{Niedermaier2019onObservability} identifizieren mangelnde Erfahrung, knappe zeitliche Kapazitäten und fehlende Ressourcen als eigenständige Herausforderung beim Aufbau von Monitoring-Konzepten in verteilten Systemen (C7). frag.jetzt wird von einem sehr kleinen Entwicklungsteam betrieben, das Observability erstmals einführt. Die Alternative aus Prometheus, Jaeger und Loki erfordert die Pflege dreier unterschiedlicher Konfigurationsformate und Storage-Backends \parencite[S.~501]{gudimetla2025transitioning}. Der Grafana-Stack reduziert diese Heterogenität, da Loki und Tempo aus demselben Ökosystem stammen und vergleichbare Konfigurationsmuster verwenden. \textcite[S.~105]{sakinala2025observability} unterstreichen, dass Integrationskomplexität zwischen verschiedenen Werkzeugen eine der häufigsten Ursachen für Verzögerungen bei der Umsetzung von Observability-Vorhaben darstellt.
 
Das zweite Argument adressiert die \emph{Korrelierbarkeit innerhalb einer einzigen Benutzeroberfläche}. Die Interviewstudie von \textcite[S.~8]{Niedermaier2019onObservability} dokumentiert das Fehlen einer zentralen, navigierbaren Gesamtsicht als wiederkehrendes Defizit (C4). \textcite[S.~16]{sundstrom2025finding} beschreiben Grafana im Kontext einer prototypischen Tracing-Implementierung als die Oberfläche, die alle Observability-Signale vereint und eine Navigation zwischen Metriken, Traces und Logs ermöglicht. \textcite[S.~86914]{usman2022survey} ordnen eine solche einheitliche Sicht, die Logs, Traces und Metriken nahezu in Echtzeit zusammenführt, als Schlüsseleigenschaft einer leistungsfähigen Observability-Plattform ein. In der alternativen Werkzeugkombination wäre Jaeger als separates Trace-Backend über die Grafana-Oberfläche zwar einbindbar \parencite[S.~17]{sundstrom2025finding}, die Verknüpfung bliebe jedoch weniger eng integriert, und die Jaeger-eigene Oberfläche würde als paralleles Werkzeug bestehen bleiben. \textcite[S.~109]{sakinala2025observability} betonen, dass kontextbezogene Dashboard-Implementierungen, die Telemetriedaten aus verschiedenen Quellen korrelieren, den Zugriff auf separate Werkzeuge während der Störungsbehandlung überflüssig machen und dadurch die Diagnosezeit verkürzen.
 
Das dritte Argument betrifft die \emph{Austauschbarkeit bei stabilem Instrumentierungskern}. OpenTelemetry entkoppelt die Instrumentierung von der Backend-Wahl (vgl.~Abschnitt~5.2). Sollte sich eine Komponente (z.B. Tempo) im Betrieb als unzureichend erweisen, ließe sich der Trace-Export im OTel Collector auf Jaeger umleiten, ohne Applikationscode anzupassen. Ebenso lässt sich die Metrikspeicherung bei wachsendem Datenvolumen von der lokalen Prometheus-TSDB auf Mimir als skalierbares Langzeit-Backend umstellen \parencite[S.~502]{gudimetla2025transitioning}. Diese Flexibilität ist für Organisationen mit begrenzter Monitoring-Erfahrung besonders relevant, deren nichtfunktionale Anforderungen sich häufig erst im laufenden Betrieb konkretisieren \parencite[S.~9--10]{Niedermaier2019onObservability}.


\subsection{Kompromisse und Einschränkungen}
 
Die Entscheidung für den Grafana-Stack geht mit bewussten Kompromissen einher. Tempo ist ein jüngeres Projekt als Jaeger, das als graduated CNCF-Projekt über einen breiteren Reifegrad verfügt \parencite[S.~15]{giamattei2024monitoring}. Die Community um Tempo ist nach Einschätzung des Autors kleiner und die öffentlich dokumentierte Lösungsbasis für Randprobleme weniger umfangreich.
 
Ebenso verzichtet die gewählte Architektur bewusst auf Mimir als Metrik-Backend. Prometheus speichert Zeitreihen lokal mit begrenzter Vorhaltezeit \parencite[S.~919]{banerjee2024comparative}. Für den aktuellen Umfang von frag.jetzt, wo moderates Telemetrievolumen, Single-Host-Betrieb und prototypische Umsetzung im Rahmen einer Bachelorarbeit gegeben sind, ist die lokale TSDB von Prometheus jedoch ausreichend. \textcite[S.~502]{gudimetla2025transitioning} beschreiben Mimir als Alternative, die horizontale Skalierung und Langzeitvorhaltung auf Objektspeicher ermöglicht. Diese Option bleibt als evolutiver Ausbaupfad bestehen.
 
Lokis labelbasierte Indexierung erfordert eine sorgfältige Auswahl der Labels, um eine unkontrollierte Vermehrung von Streams zu vermeiden \parencite{grafana_loki_docs}. Da frag.jetzt aus einer überschaubaren Zahl von Diensten besteht, lässt sich dieses Risiko durch klare Konventionen beherrschen.
 
Diese Nachteile werden durch drei Faktoren relativiert. Erstens ist die offizielle Grafana-Labs-Dokumentation für die hier genutzten Standardkonfigurationen nach Einschätzung des Autors ausreichend. Zweitens überwiegen für ein sehr kleines Team die Vorteile einer homogenen Komponentenfamilie gegenüber der breiteren Dokumentationsbasis verstreuter Einzelwerkzeuge. Drittens stellt OpenTelemetry als Instrumentierungsschicht sicher, dass ein späterer Wechsel einzelner Backend-Komponenten mit vertretbarem Aufwand möglich bleibt.
 
Damit beantwortet die gewählte Kombination aus Prometheus, Loki, Tempo und Grafana (instrumentiert über OpenTelemetry) Forschungsfrage~2. Im weiteren Verlauf bildet sie die technische Grundlage für die Konzeption der Telemetrie-Pipeline (Kapitel~6) sowie die prototypische Umsetzung (Kapitel~7).